% problem-000449 1 碱式铬酸铅沉淀
\section*{1. 碱式铬酸铅沉淀}
\textit{下列为分问。}

\subsection*{1-1}
将热的硝酸铅溶液滴⼊热的铬酸钾溶液产⽣碱式铬酸铅沉淀 $\mathrm { P b _ { 2 } ( O H ) _ { 2 } C r O _ { 4 } ] _ { \circ } }$

\subsection*{1-2}
向含氰化氢的废⽔中加⼊铁粉和 $\mathrm { K _ { 2 } C O _ { 3 } }$ 制备⻩⾎盐 $\mathrm { { K _ { 4 } F e ( C N ) _ { 6 } { \cdot } 3 H _ { 2 } O ] _ { c } } }$

\subsection*{1-3}
酸性溶液中，⻩⾎盐⽤ $\mathrm { K M n O _ { 4 } }$ 处理，被彻底氧化，产⽣ $\mathrm { N O _ { 3 } } ^ { - }$ −和 $\mathrm { { C O } _ { 2 \circ } }$

\subsection*{1-4}
在⽔中， $\mathrm { A g _ { 2 } S O _ { 4 } }$ 与单质S作⽤，沉淀变为 $\mathrm { A g } _ { 2 } \mathrm { S } ,$ ，分离，所得溶液中加碘⽔不褪⾊。

\subsection*{2-1}
实验室现有试剂：盐酸，硝酸，⼄酸，氢氧化钠，氨⽔。从中选择⼀种试剂，分别分离以下各组固体混合物（不要求复原，括号内数据是溶度积），指出溶解的固体。

\subsection*{(1)}
$\mathrm { C a C O _ { 3 } } ( 3 . 4 \times 1 0 ^ { - 9 }$ )和 $\mathrm { C a C _ { 2 } O _ { 4 } } \left( 2 . 3 { \times } 1 0 ^ { - 9 } \right)$

\subsection*{(2)}
$\mathrm { B a S O _ { 4 } } \left( 1 . 1 \times 1 0 ^ { - 1 0 } \right)$ )和 $\mathrm { B a C r O _ { 4 } } \left( 1 . 1 { \times } 1 0 ^ { - 1 0 } \right)$

\subsection*{(3)}
$\mathrm { Z n ( O H ) } _ { 2 } \left( 3 . 0 { \times } 1 0 ^ { - 1 7 } \right)$ 和 $\mathrm { N i ( O H ) } _ { 2 } \left( 5 . 5 \times 1 0 ^ { - 1 6 } \right)$

\subsection*{(4)}
$\operatorname { A g C l } \left( 1 . 8 \times 1 0 ^ { - 1 0 } \right)$ 和 $\mathrm { A g I } \left( 8 . 5 \times 1 0 ^ { - 1 7 } \right)$

\subsection*{(5)}
$\mathrm { Z n S } \left( 2 . 5 { \times } 1 0 ^ { - 2 2 } \right)$ 和 $\mathrm { H g S } \left( 1 . 6 \times 1 0 ^ { - 5 2 } \right)$

\subsection*{2-2}
在酸化的KI溶液中通⼊ $S O _ { 2 } ,$ ，观察到溶液变⻩并出现混浊(a)，继续通 $\mathrm { S O } _ { 2 }$ ，溶液变为⽆⾊(b)，写出与现象a和b相对应所发⽣反应的⽅程式。写出总反应⽅程 $\operatorname { \overline { { x } } } \mathrm { ( c ) }$ ，指出KI在反应中的作⽤。

\subsection*{2-3}
分⼦量为4000的聚⼄⼆醇有良好的⽔溶性，是⼀种缓泻剂，它不会被消化道吸收，也不会在体内转化，却能使肠道保持⽔分。

\subsection*{2-3}
-1 以下哪个结构简式代表聚⼄⼆醇？

（图：）

A

（图：）

B

（图：）

C

（图：）

D

\subsection*{2-3}
-2 聚⼄⼆醇为何能保持肠道⾥的⽔分？

\subsection*{2-3}
-3 聚⼄⼆醇可由环氧⼄烷在酸性条件下聚合⽽成，写出反应式。
